\section{Detailed results}\label{app:extra}
This appendix gives the full analysis behind Section~\ref{sec:experiments}. Comparisons follow the decision rule of Section~\ref{sec:eval}: a separately trained model must differ by more than $2\hat\sigma_{\text{seed}}$ of its base method (last row of Table~\ref{tab:training}); samplers on one model are paired comparisons and must keep their direction on the other training seeds.

\subsection{Evaluation error (RQ1)}\label{app:c1}
\begin{table}[t]
\centering
\caption{Evaluation error on MNIST (FID, lower is better). Top: base models, mean $\pm$ sd over 3 training seeds (DDIM: the 2 seeds' values), 10,000 samples. Middle: sd over 6 sampling-noise draws of the seed-0 model. Bottom: seed-0 model at NFE~100 against the number of generated samples, and the $c/n$ extrapolation FID$_\infty$ \citep{chong2020fid}.}
\label{tab:c1}
\small
\begin{tabular}{lcccc}
\toprule
 & NFE 1 & NFE 5 & NFE 100 & dopri5 \\
\midrule
I-CFM (3 seeds) & $322\pm9.28$ & $45.6\pm7.79$ & $17.0\pm5.81$ & $14.3\pm5.05$ \\
OT-CFM (3 seeds) & $212\pm9.92$ & $33.4\pm5.80$ & $16.0\pm1.88$ & $13.4\pm1.70$ \\
DDIM (2 seeds) & 293, 356 & 99.3, 106 & 12.2, 32.8 & -- \\
\midrule
sampling-noise sd & NFE 2 & NFE 10 & NFE 100 & \\
\midrule
I-CFM / OT-CFM & 0.34 / 0.58 & 0.13 / 0.28 & 0.13 / 0.05 & \\
\midrule
samples ($n$) & 2k / 5k & 10k & 20k & FID$_\infty$ \\
\midrule
I-CFM, NFE 100 & 17.3 / 14.9 & 14.5 & 14.1 & 13.7 \\
OT-CFM, NFE 100 & 17.1 / 14.5 & 13.9 & 13.4 & 13.0 \\
\bottomrule
\end{tabular}
\end{table}
% sources: results/e4/mnist_eval.json, results/e3/fid_s{1,2}.json, results/c1_reliability/fid_noise.json

\paragraph{Training seeds dominate the error.} The FID of separately trained models varies far more than the FID of one model under new sampling noise (Table~\ref{tab:c1}, Figure~\ref{fig:c1}a): at NFE~100 the seed sd is 5.8 (I-CFM) and 1.9 (OT-CFM), against a noise sd of 0.13 and 0.05, so subtracting $\sigma^2_{\text{noise}}$ from the between-seed variance changes nothing. I-CFM seed~2 is worse at every NFE (FID 23.7 at NFE~100 versus 14.5 and 12.9) although its setup and final training loss (0.169) match the other seeds; without EMA the last-epoch weights are a noisy choice. The thresholds are therefore large: $2\hat\sigma_{\text{seed}}=11.6$ (I-CFM) and 3.8 (OT-CFM) at NFE~100. % results/c1_reliability/summary.md
\paragraph{Sample size.} FID at 10,000 samples exceeds the $c/n$ extrapolation FID$_\infty$ by 0.8--0.9 for both models, so comparisons at equal $n$ are unaffected; FID from 2,000 samples is 3.6 higher and must not be mixed with 10,000-sample numbers. % results/c1_reliability/fid_noise.json
\paragraph{Method comparison.} OT-CFM has lower FID than I-CFM at 1, 2 and 5 steps on all three seeds (paired differences $-109$ to $-110$, $-49$ to $-56$ and $-1$ to $-19$). From 10 steps on the sign differs between seeds: no detectable difference. Straightness is very stable: $R=1.021$--1.022 for I-CFM and 1.010--1.011 for OT-CFM, a seed spread ten times smaller than the method gap.
\paragraph{Class coverage.} Under the two-sample test the real test set is not rejected (TV $=0.017$, $p=0.14$), and its TV is close to the floor 0.014 of \eqref{eq:tv}; the usual one-sample test, which treats $p_k$ as known, wrongly gives $p=0.037$. Generated samples reach TV 0.05--0.10 at 100 steps, with bootstrap intervals about 0.02 wide and a seed sd of 0.02, so the TV threshold is about 0.05. One-step OT-CFM covers the classes far better than one-step I-CFM (TV 0.08--0.14 versus 0.70--0.79), as \eqref{eq:postmean} predicts. I-CFM seed~2 has the worst FID but the same TV as seed~0 (0.098 versus 0.092), so its problem is image quality, not class balance. % results/c1_reliability/kmeans_c1.json

\subsection{Solvers and time grids (RQ2)}\label{app:c2}
\begin{figure}[t]
\centering
\includegraphics[width=0.85\linewidth]{figs/fig_c2_solvers.pdf}
\caption{(a) Path curvature $\kappa(t)=\E\|\tfrac{d}{dt}v(x(t),t)\|$ from 256 dopri5 trajectories (log scale) and the 8-step $\sqrt\kappa$ grid of OT-CFM (dotted lines); $G$ is the predicted error ratio \eqref{eq:sqrtk}. (b) FID (10,000 samples, lower is better) against NFE for OT-CFM seed~0: three solvers on the uniform (dashed) and $\sqrt\kappa$ (solid) grids, and Heun/RK4 whose last step does not call the network at $t=1$ (crosses). Dash-dotted: dopri5.}
\label{fig:c2}
\end{figure}
\begin{table}[t]
\centering
\caption{Inference-side samplers on the seed-0 base models: FID (10,000 samples, lower is better) at about 8, 16 and 32 network calls (exact count in parentheses when it differs), and NFE@20. The last-step variants never evaluate the network at $t=1$. Paired comparisons on one model: the sampling-noise sd is 0.05--0.6 (Table~\ref{tab:c1}).}
\label{tab:sampler}
\small
\setlength{\tabcolsep}{4pt}
\resizebox{\linewidth}{!}{%
\begin{tabular}{llcccc}
\toprule
Model & Sampler & NFE $\approx$ 8 & 16 & 32 & NFE@20 \\
\midrule
I-CFM & Euler, uniform & 32.0 & 22.4 & 18.4 & 23.7 \\
 & Euler, $\sqrt\kappa$ & 31.3 & 19.8 & 15.2 & 15.8 \\
 & Heun, uniform & 126 & 78.4 & 48.0 & $>32$ \\
 & Heun, $\sqrt\kappa$ & 61.9 & 22.1 & 13.7 & 19.0 \\
 & RK4, uniform & 125 & 65.9 & 37.1 & $>32$ \\
 & RK4, $\sqrt\kappa$ & 114 & 31.3 & 13.2 & 24.6 \\
 & Heun, last step Euler, uniform & 23.9$^{(7)}$ & 17.8$^{(15)}$ & 16.5$^{(31)}$ & 11.4 \\
 & Heun, last step Euler, $\sqrt\kappa$ & 23.8$^{(7)}$ & 16.0$^{(15)}$ & 12.9$^{(31)}$ & 10.2 \\
 & RK4, last step Heun, uniform & -- & 17.5$^{(14)}$ & 16.3$^{(30)}$ & 9.91 \\
\midrule
OT-CFM & Euler, uniform & 20.5 & 17.4 & 16.2 & 9.15 \\
 & Euler, $\sqrt\kappa$ & 19.7 & 14.3 & 12.5 & 7.88 \\
 & Heun, uniform & 134 & 83.3 & 50.4 & $>32$ \\
 & Heun, $\sqrt\kappa$ & 61.1 & 22.8 & 13.3 & 19.6 \\
 & RK4, uniform & 114 & 66.1 & 40.7 & $>32$ \\
 & RK4, $\sqrt\kappa$ & 97.0 & 31.8 & 12.8 & 24.6 \\
 & Heun, last step Euler, uniform & 18.4$^{(7)}$ & 15.5$^{(15)}$ & 15.9$^{(31)}$ & 6.60 \\
 & Heun, last step Euler, $\sqrt\kappa$ & 19.0$^{(7)}$ & 15.6$^{(15)}$ & 12.4$^{(31)}$ & 6.67 \\
 & RK4, last step Heun, uniform & -- & 16.4$^{(14)}$ & 16.5$^{(30)}$ & 7.73 \\
\bottomrule
\end{tabular}}
\end{table}
% sources: results/e4/mnist_eval.json, results/c2_solvers/{icfm,otcfm}_s0.json, results/c2_solvers/t1_{icfm,otcfm}_s0.json

\paragraph{Curvature.} For both models $\kappa(t)$ rises about a hundredfold for $t>0.95$, and 65--69\% of $\int\kappa$ lies in $t>0.9$ (Figure~\ref{fig:c2}a). The predicted error ratios are $G=0.50$ (OT-CFM) and 0.54 (I-CFM). Early in the path OT-CFM bends less than I-CFM, as its straighter pairs suggest. % results/c2_solvers/{otcfm,icfm}_s0.json kappa
\paragraph{The $\sqrt\kappa$ grid.} With Euler it lowers FID at 16 and 32 steps from 17.4 to 14.3 and from 16.2 to 12.5 (OT-CFM), and from 22.4 to 19.8 and from 18.4 to 15.2 (I-CFM); at 2--4 steps it is equal or slightly worse (Table~\ref{tab:sampler}). These are paired comparisons on one model, so a 3-point difference is 10--60 sampling-noise sd. OT-CFM with Euler on the $\sqrt\kappa$ grid comes within 10\% of the dopri5 FID with about 30 NFE, against about 270 for dopri5.
\paragraph{The $t=1$ effect.} On the uniform grid Heun and RK4 are far worse than Euler at every NFE (OT-CFM at 32 NFE: 50.4 and 40.7 versus 16.2). Heun's corrector and RK4's last stage evaluate the network at $t=1$, which Euler never does. There, by Proposition~\ref{prop:onestep} taken at $t=1$, the ideal output is $u_1(x)=x-\E[x_0\mid x_1=x]=x$ for independent pairs (and nearly so for OT-CFM, whose minibatch pairs are close to independent in 784 dimensions), but our networks miss it by 43\% (RMS relative error) on real images and by about 85\% on generated ones; the RMS error grows from 0.36 at $t=0.9$ to 0.57 at $t=0.999$. Replacing only the last step (Euler for Heun; Heun with its end evaluation at $1-h/2$ for RK4) removes the failure: at 31 and 30 NFE OT-CFM reaches 15.9 and 16.5 instead of 50.4 and 40.7, and I-CFM 16.5 and 16.3 instead of 48.0 and 37.1. These paired differences of 21--35 FID points exceed the sampling-noise sd more than 50 times. The $\sqrt\kappa$ grid partly hides the effect because its last step is so short that the bad evaluation carries little weight. % results/c2_solvers/t1_{otcfm,icfm}_s0.json
\paragraph{Steps to FID 20.} The best inference-only sampler reduces NFE@20 from 9.2 to 6.6 for OT-CFM (Heun* on the $\sqrt\kappa$ grid) and from 23.7 to 9.9 for I-CFM (RK4 with a Heun last step), at no training cost. With many steps the gains shrink: at 31--32 NFE, Heun* on the $\sqrt\kappa$ grid gives 12.4 (OT-CFM) and 12.9 (I-CFM), Euler on the same grid 12.5 and 15.2.
\paragraph{Endpoint error does not predict FID.} Choosing, for each NFE, the grid with the smallest endpoint error against dopri5 on 256 calibration paths always selects steps concentrated near the noise ($s\approx0.2$--0.4 in the shift family) and makes FID worse than the uniform grid in all ten Euler cases, e.g.\ from 76.0 to 108.8 at 2 steps for OT-CFM; across pairs of grids, the endpoint-error ranking agrees with the FID ranking in only 124 of 204 cases (61\%). Endpoint error is dominated by early errors that decide which digit is drawn, whereas FID responds to the fine detail settled near $t=1$, where $\kappa$ peaks. Full tables are in \texttt{results/c2\_solvers/summary.md}.

\begin{table}[h]
\centering
\caption{The inference-side conclusions of Section~\ref{sec:c2} on all three training seeds (FID, 10,000 samples, lower is better; ``a $\to$ b'' is a paired comparison on one model). Columns 3--5: Euler, uniform $\to$ $\sqrt\kappa$ grid at 4, 16 and 32 steps. Column 6: Heun on the uniform grid at 32 NFE $\to$ Heun whose last step is Euler (31 NFE). Column 7: the latter on the $\sqrt\kappa$ grid. Last column: RMS error of $v(x,1)$ against its ideal value $x$, in \% of $\mathrm{RMS}(x)$, on real / generated images.}
\label{tab:seeds}
\small
\setlength{\tabcolsep}{3pt}
\resizebox{\linewidth}{!}{%
\begin{tabular}{llcccccc}
\toprule
Model & Seed & Euler 4 & Euler 16 & Euler 32 & Heun 32 $\to$ Heun* 31 & Heun* 31, $\sqrt\kappa$ & $t=1$ error (\%) \\
\midrule
I-CFM & 0 & 56.6 $\to$ 60.6 & 22.4 $\to$ 19.8 & 18.4 $\to$ 15.2 & 48.0 $\to$ 16.5 & 12.9 & 42 / 87 \\
 & 1 & 47.8 $\to$ 54.2 & 18.6 $\to$ 16.5 & 15.9 $\to$ 12.8 & 58.1 $\to$ 14.4 & 11.6 & 47 / 83 \\
 & 2 & 64.1 $\to$ 69.3 & 31.2 $\to$ 29.8 & 27.5 $\to$ 24.4 & 52.8 $\to$ 24.5 & 21.9 & 37 / 91 \\
\midrule
OT-CFM & 0 & 32.3 $\to$ 32.8 & 17.4 $\to$ 14.3 & 16.2 $\to$ 12.5 & 50.4 $\to$ 15.9 & 12.4 & 43 / 84 \\
 & 1 & 43.8 $\to$ 44.5 & 21.2 $\to$ 20.7 & 18.8 $\to$ 17.2 & 50.4 $\to$ 17.1 & 15.5 & 55 / 83 \\
 & 2 & 42.6 $\to$ 43.8 & 22.5 $\to$ 20.5 & 20.2 $\to$ 17.0 & 57.6 $\to$ 18.9 & 15.4 & 34 / 83 \\
\bottomrule
\end{tabular}}
\end{table}
% source: results/c2_solvers/{icfm,otcfm}_s{0,1,2}.json and t1_*.json (code/c2_solvers/sample_grid.py --core, t1_check.py)

\paragraph{Other training seeds.} Table~\ref{tab:seeds} repeats the main comparisons on seeds 1 and 2. Every direction found on seed~0 holds on all six models. The $\sqrt\kappa$ grid lowers Euler FID at 16 steps (by 0.5--3.1) and at 32 steps (by 1.6--3.7), and is slightly worse at 4 steps (by 0.5--6.4). On the uniform grid Heun is 1.9--3.7 times worse than Euler at 32 NFE, and moving its last network call away from $t=1$ removes that loss (to 14.4--24.5). The network misses its ideal value $v(x,1)=x$ by 34--55\% on real images and by 83--91\% on generated ones on every seed. The size of the $\sqrt\kappa$ gain varies between models (OT-CFM seed~1 gains 0.5 at 16 steps, seed~0 gains 3.1), so we report it per model rather than as one number. % results/c2_solvers/{icfm,otcfm}_s{1,2}.json, t1_*_s{1,2}.json

\subsection{Training-side variants (RQ3)}\label{app:c35}
\begin{table}[t]
\centering
\caption{Pairing cost $\E\|x_0-x_1\|^2$ saved relative to random pairs on MNIST (random-pair cost 1,510), as a share of that cost and of the Gelbrich bound $2\operatorname{tr}\Sigma^{1/2}=386$, and the CPU time of one exact OT solve (POT, 4 threads).}
\label{tab:pairs}
\small
\setlength{\tabcolsep}{4.5pt}
\adjustbox{max width=\linewidth}{%
\begin{tabular}{lcccccccc}
\toprule
OT batch & 16 & 64 & 128 & 256 & 512 & 1,024 & 2,048 & 4,096 \\
\midrule
saved / random cost & 2.9\% & 4.0\% & 4.5\% & 4.8\% & 5.2\% & 5.6\% & 5.8\% & 6.1\% \\
saved / bound & 11\% & 16\% & 18\% & 19\% & 21\% & 22\% & 23\% & 24\% \\
OT solve (s) & 0.0003 & 0.00054 & 0.0019 & 0.0069 & 0.038 & 0.15 & 1.2 & 6.2 \\
\midrule
\multicolumn{9}{l}{Reflow pairs from I-CFM seed 0 (60,000 pairs): cost 1319, saved 12.7\% of the random cost, 50\% of the bound.} \\
\multicolumn{9}{l}{Reflow pairs from OT-CFM seed 0 (60,000 pairs): cost 1316, saved 12.9\% of the random cost, 50\% of the bound.} \\
\bottomrule
\end{tabular}}
\end{table}
% sources: results/c3_otbatch/pairs.json, results/c5_reflow/pairs_{icfm,otcfm}_s0.json

\paragraph{Minibatch OT saves little at any batch size.} Exact OT within larger batches saves only slowly more pairing cost: 4.5\% of the random cost at batch 128, 5.8\% at 2,048 and 6.1\% at 4,096, i.e.\ 18--24\% of the bound \eqref{eq:bound} (Table~\ref{tab:pairs}). Each doubling saves about 0.4\% more; extrapolating the logarithmic trend, 10\% would need batches of about $2.8\times10^6$ samples, beyond the 60,000 training images, while one OT solve already takes 1.2\,s at batch 2,048 and 6.2\,s at 4,096. % results/c3_otbatch/pairs.json
\paragraph{OT batch 512.} The one-step FID is 196, lower than all three base seeds (205--223). Against seed~0 the gain (27.6) exceeds the threshold of 19.8, but seed~0 is the worst base seed at one step, and against the three-seed mean (212) the gain is 16.3, within the threshold; we therefore do not count it as a detectable improvement. The paths are marginally straighter ($R=1.0098$ versus 1.010--1.011). From five steps on the variant is 3--5 points worse at every NFE, within the threshold except with dopri5 (16.7 versus 11.5; threshold 3.4). OT solving took 279\,s of 1.56\,h. % results/e3/fid_s0_otb512.json
\paragraph{OT batch 2,048.} The one-step FID (208) lies inside the base-seed range. From two steps on the variant is 3--11 points worse than seed~0 (82.9 versus 76.0 at 2 steps, 17.4 versus 13.9 at 100), beyond the threshold only with dopri5 (15.6 versus 11.5) and within it at every NFE against the three-seed mean. Its paths are the straightest of all OT-CFM models ($R=1.0085$) and its final loss the lowest (0.100 versus 0.118), but neither shows in FID; the lower loss is expected from Proposition~\ref{prop:cfm}, since a better coupling lowers the constant $C$, and says nothing about the fit. The run was interrupted by a power loss after epoch 21 and resumed from the epoch-20 checkpoint; its time, 2.37\,h, adds both parts, of which 1,693\,s (20\%) is OT solving. % results/e3/fid_s0_otb2048.json; results/c3_otbatch/logs/train_otcfm_s0_otb2048.log (3021 s + 5522 s)
\paragraph{Logit-normal training times.} The variant has FID 327, 103, 30.0, 18.3, 14.1, 12.0 and 12.6 at 1, 2, 5, 10, 20, 50 and 100 steps. By the strict rule only the two-step gain over seed~0 (29.4; threshold 10.2) is detectable; at five steps the gain (15.3) is just below the threshold (15.6). The variant is, however, better than all three base seeds at every NFE from 2 to 100. If it were merely a fourth draw from the seed distribution, it would rank first at any one NFE with probability $1/4$; the ranks at different NFEs are strongly correlated, so this is suggestive rather than a test. NFE@20 falls from 23.7 to 9.1, as low as the best inference-side sampler achieves on the base model (9.9). Two observations are consistent with less training near $t=1$, where the logit-normal density vanishes: its dopri5 FID (14.2, at only 111 NFE) is worse than its own 50-step Euler FID (12.0), and dopri5, unlike Euler, evaluates the network at $t=1$. % results/e3/fid_s0_tlogitnormal.json
\paragraph{$\kappa$-weighted training times.} Putting half of the training times where the paths bend gives nothing beyond the threshold: FID is 10.0, 8.1 and 1.9 lower than seed~0 at 2, 5 and 10 steps, and 2.1--3.8 higher at 20--100 steps and 5.3 higher with dopri5. The paths become the curviest of all models ($R=1.024$), and dopri5 needs 345 NFE instead of 280. The sampler should spend its steps near $t=1$, but training should not spend its samples there. % results/e3/fid_s0_tkappa.json
\paragraph{Reflow from I-CFM.} We integrated I-CFM seed~0 with dopri5 from 60,000 fresh noise vectors (287 NFE on average) and fine-tuned the same network on these pairs. The reflow pairs cost 1,319, a saving of 12.7\% of the random cost and half of the bound \eqref{eq:bound}, twice what minibatch OT achieves even at batch 4,096 (Table~\ref{tab:pairs}). The training loss falls from 0.168 to 0.004: with a deterministic coupling the constant $C$ of Proposition~\ref{prop:cfm} is close to zero. The paths become almost straight ($R=1.0008$ versus 1.021), and few-step FID improves far beyond the threshold: 332 to 64.6 at one step, 132 to 30.9 at two and 45.3 to 23.9 at five. One reflow step gives the same digit as 100 base steps from the same noise (Figure~\ref{fig:samples}). With more steps the model plateaus: 22.0 at 20 steps and 20.5 at 100, against 20.7 and 14.5 for the base model, and 19.0 instead of 12.1 with dopri5. These losses are below the I-CFM threshold but go in the same direction at every NFE from 20 on. With an exact field the reflowed model would reproduce the base model's dopri5 samples, so 12.1 is the best it can reach (Appendix~\ref{app:d7}); we attribute the remaining gap of about 7 points to the error of the second regression, which adds to that of the first, but have not separated this from the effect of fine-tuning rather than training from scratch. The lower ceiling is not a loss of class balance: cluster TV is 0.066--0.069 at every NFE from 1 to 100 (bootstrap interval 0.058--0.078 at one step), against 0.79, 0.25 and 0.13 for the base model at 1, 2 and 5 steps and 0.085--0.092 at 100. % results/e3/fid_s0_reflow.json, results/c5_reflow/pairs_icfm_s0.json, results/c5_reflow/kmeans_reflow.json
\paragraph{Reflow from OT-CFM.} The same procedure on OT-CFM seed~0 (277 NFE per pair) gives pairs of cost 1,316 (12.9\% saved) and FID 61.2, 28.9, 23.8, 22.6, 22.1, 21.7 and 21.5 at 1 to 100 steps, 21.1 with dopri5, and $R=1.0007$. The one- and two-step gains are again far beyond the threshold. Because OT-CFM's seed spread is smaller, the lower ceiling is now detectable: the model is worse than its base by 5.1, 6.2 and 7.6 at 20, 50 and 100 steps and by 9.6 with dopri5 (thresholds 4.7, 3.9, 3.8 and 3.4). After reflow the two methods are almost identical (differences of 3.4 at one step and at most 2.1 elsewhere), so OT-CFM's base advantage of about 110 one-step FID points does not survive reflow. We did not compute cluster TV for this model. % results/e3/fid_s0_reflow.json, results/c5_reflow/pairs_otcfm_s0.json
\paragraph{Cost.} Reflow costs 2.7 (I-CFM) and 2.8 (OT-CFM) times the base training time: pair generation 2.05 and 1.84\,h plus retraining 1.08 and 1.56\,h, against 1.17 and 1.20\,h for the base model, all throttled (Appendix~\ref{app:setup}). % results/e2/*_train.json, results/c5_reflow/pairs_*.json

\subsection{Samplers on the reflow models, and straightness (RQ4)}\label{app:c6}
\paragraph{Samplers stack with reflow only partly.} On both reflow models the $\sqrt\kappa$ grid lowers Euler FID from 4 steps on by 0.5--2.4 (I-CFM reflow: 21.7 to 19.4 at 32 steps; OT-CFM reflow: 24.7 to 22.3 at 4) and is worse at 2 steps for I-CFM (32.4 versus 30.9). On the uniform grid Heun and RK4 still fail although the paths are nearly straight (Heun at 32 NFE: 46.2 and 34.6, against 21.7 and 21.9 for Euler), so the failure is not caused by curvature. We did not run the $t=1$ diagnostic or Heun* on the reflow models; for a deterministic coupling the ideal $u_1(x)=x-\E[x_0\mid x_1=x]$ is no longer $x$, so the argument of Section~\ref{sec:c2} does not carry over directly. % results/c2_solvers/{icfm,otcfm}_s0_reflow.json
\paragraph{The lower envelope.} Taking every curve of Figure~\ref{fig:frontier}, the envelope is formed by reflow at 1--4 NFE (61.2, 28.9 and, with the $\sqrt\kappa$ grid, 22.3 at 4) and, from 7 NFE on, alternately by OT-CFM with Heun* or Euler on the $\sqrt\kappa$ grid and by the logit-normal I-CFM model with plain Euler (19.0 at 7, 18.3 at 10, 14.3 at 16, 14.1 at 20, 12.4 at 31, 12.0 at 50). At similar NFE these curves are within about one FID point of each other, well below the seed threshold, so the envelope does not rank them.
\paragraph{Straightness and one-step FID.} Across the twelve models (six base models, four OT-batch and time-distribution variants, two reflow models), the Spearman correlation between $R$ and one-step FID is 0.92 (0.85 without the reflow models). The order is driven by the gaps between groups (reflow $<$ OT-CFM $<$ I-CFM); within OT-CFM (five models) it is 0.3 and within I-CFM (five models) 0.5. The models are not independent (variants share seed~0), so we report these numbers descriptively.

\begin{figure}[h]
\centering
\includegraphics[width=\linewidth]{figs/fig_2d_paths.pdf}
\caption{2D illustration on the moons data (training seed 0). Curves: 100-step Euler trajectories from Gaussian noise (dots). Crosses: where a single Euler step lands. I-CFM's one-step samples collapse to the data mean (Appendix~\ref{app:d3}); OT-CFM's straight paths make one step nearly exact. In 2D minibatch OT saves 87--93\% of the pairing cost, on MNIST 4.5\% (Table~\ref{tab:pairs}).}
\label{fig:2d}
\end{figure}

\begin{table}[t]
\centering
\caption{Published unconditional CIFAR-10 results for context (FID on 50,000 samples, numbers as reported). They use far larger networks and training budgets and a different dataset, so they are not comparable with our MNIST numbers; they show the same pattern of NFE versus quality.}
\label{tab:literature}
\small
\begin{tabular}{lcc}
\toprule
Method & NFE & FID \\
\midrule
DDPM \citep{ho2020ddpm} & 1,000 & 3.17 \\
DDIM \citep{song2021ddim} & 10 / 100 & 13.36 / 4.16 \\
EDM, Heun with tuned grid \citep{karras2022edm} & 35 & 1.97 \\
I-CFM / OT-CFM, dopri5 \citep{tong2024otcfm} & 146 / 134 & 3.66 / 3.58 \\
1-rectified flow, one Euler step \citep{liu2023rf} & 1 & 378 \\
2-rectified flow (reflow), one Euler step \citep{liu2023rf} & 1 & 12.21 \\
Consistency distillation \citep{song2023cm} & 1 / 2 & 3.55 / 2.93 \\
\bottomrule
\end{tabular}
\end{table}
% sources: papers/10 (p.1), 11 (Table 1, eta = 0), 16 (Table 2), 23 (Table 5), 21 (Table 1), 31 (Table 1)

